% Front matter: edition notice (PUB-REF-0001 AN-20, the house verso; the
% slogan under the edition notice per the operator directive of 2026-10-17)
% and How to Use This Book (AN-21, F-MAC: Part I read in order, Part II used
% as an encyclopedia, the Summary, the conventions; the 1991 Macintosh
% User's Guide preface, apple-macintosh.md Sec 1, 4).  All boilerplate newly
% written (VO-12).  The colophon is pubs/common/colophon.tex (AN-26).
% GATES: test-pubs-build test-pubs-flair
\Gates{test-pubs-build, test-pubs-flair}
\begin{TitleVerso}
\IsThisGoodForTheCompany
{\imLegalFont
\textbf{The following paragraph does not apply to any country where such
provisions are inconsistent with local law:} INITECH SYSTEMS CORPORATION
PROVIDES THIS PUBLICATION ``AS IS'' WITHOUT WARRANTY OF ANY KIND, EITHER
EXPRESS OR IMPLIED, INCLUDING, BUT NOT LIMITED TO, THE IMPLIED WARRANTIES OF
MERCHANTABILITY AND FITNESS FOR A PARTICULAR PURPOSE.  Some states do not
allow the exclusion of implied warranties, so the above exclusion may not
apply to you.\par
This publication may contain technical inaccuracies or typographical errors.
The information in it is revised from time to time, and the revisions are
published in new editions.  Initech Systems Corporation may change the
products and programs described in this publication at any time without
notice.\par
Requests for copies of Initech publications and for technical information
about Initech products should be made to your authorized Initech dealer or
your Initech representative.\par}

{\footnotesize
A Reader's Comment Form is provided at the back of this publication.  If the
form has been removed, address comments to: \InitechAddressLine.
\draft{PUB-REF-0001 O-5 and D-12: the address and permit number are to be
checked against real businesses before print}
Initech may use or distribute the information you supply in any way it
believes appropriate without incurring any obligation to you.\par
The software described in this publication is furnished under the Initech
Program License Agreement and may be copied only as that agreement
permits.\par}
\end{TitleVerso}

\FrontChapter{How to Use This Book}
\index{book, how to use}
This book tells you how to work with FLAIR\tm{}, the graphical desktop of
InitechOS\tm{}.  FLAIR shows your disks, folders and files as small pictures
on the screen, runs your applications in windows, and lets you give most
commands by pointing with the mouse and choosing from menus.

The book has two parts, a summary, four appendixes, a glossary and an index.

\section{Part I: Learning FLAIR}
Part I is a course.  Read its chapters in order, at the computer, and do
each task as you read about it.  When you have finished Part I you will be
able to start and stop the computer, use the mouse and the menus, work with
windows, organize your files in folders, open and quit applications, and
throw away what you no longer need.

\section{Part II: FLAIR Reference}
Part II is for looking things up.  Use it as you would use an encyclopedia:
turn to the chapter on the subject you want, or find the subject in the
index.  Chapter~8 describes every command in every FLAIR menu, menu by
menu, in the order the menus appear in the menu bar.  Chapter~11 lists
every alert and dialog box FLAIR displays, and Chapter~14 tells you how FLAIR
and the InitechDOS command prompt work together.

\section{The Summary}
The Summary, after Part II, collects the most common tasks in short tables:
what you want to do on the left, the steps that do it on the right.  Once
you have worked through Part I, the Summary is often all you need.

\section{The Appendixes}
Appendix~A lists every file on the InitechOS data disk, the disk FLAIR
shows as the INITECH icon.  Appendix~B describes what your computer needs
to run FLAIR, Appendix~C the keyboard and the characters you can type, and
Appendix~D the other books and the Help Desk.

\section{Other Books}
\SideHead{Getting started}\doctitle{Getting Started with FLAIR}
(1031-0000-11) tells you how to set up your computer and start FLAIR for
the first time, and what to do if FLAIR does not start.

\SideHead{Accessories}The \doctitle{FLAIR Accessories Guide} (1031-0000-12)
describes the Calculator, InitechText\tm{}, InitechCalc\tm{},
InitechPaint\tm{} and the other small applications that come with FLAIR.

\SideHead{Quick reference}The \doctitle{FLAIR Quick Reference}
(1031-0000-20) is a card that shows every menu and every keyboard
equivalent.

\SideHead{InitechDOS}The \doctitle{InitechDOS User's Guide} (1024-3300-11)
and the \doctitle{InitechDOS Reference} (1024-3300-10) describe the command
prompt you reach when you leave FLAIR.

Initech 123, InitechWord, InitechBase and the TPS Report Generator each come
with books of their own.

\section{Conventions Used in This Book}
\SideHead{Steps}Tasks are written as numbered steps.  The step itself is in
bold type; the plain text under it tells you what happens when you do it.

\SideHead{Menu commands}When you are told to choose a command from a menu,
the command and the menu are in bold type: ``Choose \menuitem{New Folder}
from the \menuname{File} menu.''  The command is spelled exactly as it
appears in the menu.

\SideHead{Keys}The names of keys are shown like this: \key{Ctrl},
\key{Enter}, \key{Esc}.  \key{Ctrl-N} means: hold down \key{Ctrl}, press
\key{N}, and release both keys.

\SideHead{Keyboard equivalents}Many menu commands can also be given from
the keyboard.  A menu shows the keyboard equivalent at the right of the
command, as a caret and a letter: \CmdKey{N} beside \menuitem{New Folder}
means that you can hold down \key{Ctrl} and press \key{N} instead of
choosing the command.

\SideHead{Notes and warnings}\textit{Note:} marks a helpful remark.  A
shaded box headed \textbf{WARNING} tells you how to avoid losing
information.

\SideHead{Screens}The figures are reproductions of the FLAIR screen exactly
as it appeared when the task was carried out, with the clock set to 9:30
in the morning.  Your screen shows the time of day.
